% =====================================================================
%  §7 A∞-structure for the derived category of coherent sheaves
%  （原文 p48 第 25 行起，至 p52 第 17 行命题 10 止）
% =====================================================================
\section{凝聚层导出范畴的 \texorpdfstring{$\Ainf$}{A∞}-结构}

\subsection{刚性解析空间}

对本节的读者来说，熟悉非阿基米德分析（non-archimedean analysis）的一些基本事实
（见 \cite{BGR}）会有帮助（尽管并非必需）。对任意带整仿射结构的光滑流形 $Y$，
我们将构造 $Y$ 上一个由 $\C_{\eps}$-代数构成的层 $\OO_Y$。该层的茎 $\OO_{Y,y}$
都是诺特代数，因而可以定义 $\OO_Y$-模的凝聚层这一概念。若 $Y=\R^n/\Z^n$ 是带标准
整仿射结构的环面，则凝聚 $\OO_Y$-模的范畴将等价于（由非阿基米德版本的 GAGA）
$\C_{\eps}$ 上阿贝尔簇的凝聚层范畴。

我们从局部情形开始。记 $v:\C_{\eps}\to\R\cup\{+\infty\}$ 为由
\[
  v\Bigl(\sum_{\lam_1<\lam_2<\cdots}c_i e^{-\lam_i/\eps}\Bigr)=\lam_1\quad(c_1\neq 0),
  \qquad v(0)=+\infty
\]
定义的（非离散）赋值。

\begin{dfn}
设 $U\subset\R^n$ 是标准向量空间 $\R^n$ 的开子集。定义 $\OO_{\R^n}(U)$ 为 $\C_{\eps}$ 上
由形式 Laurent 级数
\[
  f=\sum_{k_1,\dots,k_n\in\Z^n}a_{k_1\dots k_n}z_1^{k_1}\cdots z_n^{k_n}
\]
构成的向量空间，其中 $z_1,\dots,z_n$ 是形式变量，$a_{k_1\dots k_n}\in\C_{\eps}$，
并且对任意 $(y_1,\dots,y_n)\in U$ 都有
\[
  \lim_{\sum_i|k_i|\to\infty}\Bigl(v(a_{k_1\dots k_n})+\sum_i k_i y_i\Bigr)=+\infty .
\]
\end{dfn}

由定义可知，若 $f\in\OO_{\R^n}(U)$ 且 $(z_1,\dots,z_n)\in(\C_{\eps}^*)^n$，
则只要 $(v(z_1),\dots,v(z_n))\in U$，级数
$\sum_{k_1,\dots,k_n}a_{k_1\dots k_n}z_1^{k_1}\cdots z_n^{k_n}$
就在 adic 拓扑下收敛。

我们如下引入群 $GL(n,\Z)\ltimes\R^n$ 在 $(\R^n,\OO_{\R^n})$ 上的作用：
\begin{enumerate}[label=\alph*)]
  \item $GL(n,\Z)$ 同时作用在坐标的线性变换与级数中指标 $(k_1,\dots,k_n)$ 的线性变换上；
  \item 平移 $(t_1,\dots,t_n)\in\R^n$ 作用在坐标 $(y_1,\dots,y_n)$ 上，就是把
        $(y_1,\dots,y_n)$ 变为 $(y_1+t_1,\dots,y_n+t_n)$；它作用在级数上则是重新标度系数
\[
  \sum_{k_1,\dots,k_n}a_{k_1\dots k_n}z_1^{k_1}\cdots z_n^{k_n}
  \mapsto
  \sum_{k_1,\dots,k_n}\bigl(a_{k_1\dots k_n}e^{-\sum_i t_i k_i/\eps}\bigr)
  z_1^{k_1}\cdots z_n^{k_n}.
\]
\end{enumerate}

利用这个作用，我们对任意带整仿射结构的光滑流形 $Y$ 定义层 $\OO_Y$。

我们断言，存在一个典范地联系于 $Y$ 的、定义在 $\C_{\eps}$ 上的刚性解析空间
（rigid analytic space）$Y^{an}$。构造如下。考虑 $Y$ 的一个开覆盖 $\{U_i\}$，使得所有
非空交 $U_{i_1i_2\dots i_k}:=U_{i_1}\cap\cdots\cap U_{i_k}$ 在某个局部仿射坐标下
都是面具有有理斜率的凸多面体。每个 $U_{i_1i_2\dots i_k}$ 都可以等同于有限多个
半空间的交，使得它们在映射 $v^n:(\C_{\eps}^*)^n\to\R^n$ 下的原像都是形如
\[
  \bigl\{(z_1,\dots,z_n)\ \bigm|\ v(z_1^{k_1}\cdots z_n^{k_n})\ge C\bigr\}
\]
的集合，其中 $C>0$ 为有理数。由 Tate 的工作可知，这样的不等式组定义一个仿射型区域
（affinoid domain），即 $\C_{\eps}$ 上刚性解析空间的局部模型。

\begin{dfn}
定义 $Y^{an}$ 为借助 $GL(n,\Z)\ltimes\R^n$ 的作用、把局部数据 $(U_i,\OO_{U_i})$
粘合起来所得的 $\C_{\eps}$ 上刚性解析空间。
\end{dfn}

容易看出，$Y^{an}$ 是典范地定义的，并且 $Y^{an}$ 上凝聚解析层（解析几何意义下）的范畴
$\Coh(Y^{an})$ 等价于凝聚 $\OO_Y$-模（即局部有限生成的 $\OO_Y$-模）的范畴。

对 $\C_{\eps}$ 上每个代数簇 $\YY$，都可以典范地联系一个刚性解析空间 $\YY^{an}$。
若 $\YY$ 射影，则范畴 $\Coh(\YY^{an})$ 等价于 $\YY$ 上代数凝聚层的范畴
$\Coh(\YY)$（GAGA 定理）。

设 $Y=\R^n/\Lambda$ 是由 $\Z^n\subset\R^n$ 诱导的标准整仿射结构所带的 $n$ 维环面，
且 $\Lambda$ 与 $\Z^n$ 可公度。下述结果可由 \cite{Mum} 推出。

\begin{prp}
在上述记号下，有 $Y^{an}\simeq\YY^{an}$，其中 $\YY$ 是 $\C_{\eps}$ 上的阿贝尔簇。
\end{prp}

现在回到阿贝尔簇情形下度量坍缩的图景。由于坍缩是由格的重新标度定义的（见 §2），
可以证明 $\YY$ 同构于 $\C_{\eps}$ 上原来的阿贝尔簇。因此，在阿贝尔簇的情形，
我们对坍缩有两种等价的描述：一种用黎曼几何，一种用非阿基米德解析几何。

\begin{rmk}
对于带奇异纤维的坍缩，如上构造的刚性解析空间 $Y^{an}$ 似乎是「错误」的。首先，
它不紧致，因为 $Y$ 不紧致。但还有一个更根本的问题：$Y^{an}$ 似乎不能嵌入到与射影
代数簇相联系的紧解析空间中去。有若干迹象表明，存在另一个代数层 $\OO'_Y$，它
（在 $Y$ 上局部地）同构于 $\OO_Y$，而与 $(Y,\OO'_Y)$ 相联系的刚性解析空间
$(Y^{an})'$ 容许一个代数紧化。一般地，作为 $\OO_Y$ 的扭曲版本的层 $\OO'_Y$ 由第一
非阿贝尔上同调 $H^1(Y,\underline{\Aut}(\OO_Y))$ 分类，其中 $\underline{\Aut}(\OO_Y)$
是 $\OO_Y$ 的自同构群层。因此，对于不是阿贝尔簇的 Calabi--Yau 流形的镜像对称，
我们期待一个新的要素，即上同类 $[\OO'_Y]$。
\end{rmk}

\begin{yizhu}[扭曲层与第一非阿贝尔上同调]
层 $\OO_Y$ 的扭曲版本由 $H^1(Y,\underline{\Aut}(\OO_Y))$ 分类，这一点与复几何中线丛由
$H^1(Y,\OO_Y^*)$ 分类是同一机制：把 $\OO'_Y$ 在覆盖 $\{U_i\}$ 上的粘合数据看成取值于
自同构群层的 $1$-上闭链即可。因此，除了度量的示性类 $[g]$ 之外，解析图景还需要
指定这个新的上同类 $[\OO'_Y]$。
\end{yizhu}

\begin{yizhu}[刚性解析空间与赋值]
刚性解析空间（rigid analytic space）是 Tate 建立的、复解析空间的非阿基米德
类比：其局部模型是仿射型区域，即由有限多个不等式 $v(f_i)\ge c_i$ 在
$(\C_{\eps}^*)^n$ 中截出的区域，其上的函数由收敛的 Laurent 级数给出（系统叙述见
\cite{BGR}）。本文的 $Y^{an}$ 就是把 $Y$ 的每个整仿射坐标卡 $U\subset\R^n$ 实现为
$(\C_{\eps}^*)^n$ 中由 $v(z_i)=y_i$ 给出的仿射型区域，再沿
$GL(n,\Z)\ltimes\R^n$ 的作用粘合而得。Berkovich 的观点则是在此之上补充所有可乘
赋值点，得到局部紧的解析空间（见 \cite{Be}）。
\end{yizhu}

\begin{yizhu}[命题 9 的含义]
命题 9 的实质是 Mumford 的构造：在本文的设定下，$\C_{\eps}$ 上具有完全退化约化的
阿贝尔簇可以写成 $(\C_{\eps}^*)^n$ 用某个格 $\Lambda$ 所作的商，而赋值映射
$v^n:(\C_{\eps}^*)^n\to\R^n$ 恰好把这个商投影到实环面 $\R^n/\Lambda$ 上，
于是实环面上的整仿射结构完整地记录了该阿贝尔簇的解析模型（见 \cite{Mum}、
\cite{BL}）。这也解释了为什么在阿贝尔簇的情形，度量坍缩的两种描述——黎曼几何的
与非阿基米德解析的——会给出同一个对象。注意右端的 $\YY$ 是 $\C_{\eps}$ 上的阿贝尔簇，
与作为实环面的 $Y$ 不是同一个对象（原文用花体 $\YY$ 与 $Y$ 区分），
$Y$ 所记录的只是 $\YY$ 的整仿射数据。
\end{yizhu}

\subsection{凝聚层上的 \texorpdfstring{$\Ainf$}{A∞}-结构}

$Y$ 上有一个阿贝尔群层 $Af_Y$，它由整斜率的局部仿射函数组成（这样的函数局部地
写成 $l=c+\sum_{1\le i\le n}m_i y_i$，其中 $m_i\in\Z$，$c\in\R$）。
存在层的态射 $\exp:Af_Y\to\OO_Y^*$，它把 $l$ 送到
\[
  \exp(l):=e^{-c/\eps}\prod_{1\le i\le n}z_i^{m_i}.
\]


设 $(Y,g)$ 是 AK-流形（见 §3.2）。我们要定义度量的一个示性类 $[g]$，它将是复几何中
凯勒形式的上同调类的类比。设 $Af_Y\otimes\R$ 为所有实值局部仿射函数构成的层。
对于凸集覆盖 $(U_i)_{i\in I}$，可以选取光滑函数 $K_i$，使得 $g|_{U_i}=\p^2K_i$。
于是 $K_i-K_j\in(Af_Y\otimes\R)(U_i\cap U_j)$，它定义了一个 $1$-上闭链，其上同调类
记为 $[g]\in H^1(Y,Af_Y\otimes\R)$。若对偶仿射结构（见 §3）是整的，我们就得到
子群 $H^1(Y,Af_Y)/\mathrm{torsion}\subset H^1(Y,Af_Y\otimes\R)$ 中的一个类 $[g]$。
这样的类称为整的。此时 $\exp([g])\in H^1(Y^{an},\OO_Y^*)$ 是 $Y^{an}$ 上某个线丛的
第一 Chern 类。与凯勒几何类比，我们期待这个线丛是丰富的。当 $(Y,g)$ 是平坦环面时，
丰富性可以直接证明（见 \cite{BL}）。

\begin{yizhu}[AK-流形]
AK-流形是 §3.2 引入的概念，「AK」取自 affine and Kähler（仿射凯勒）：
Monge--Ampère 流形是三元组 $(Y,g,\nabla)$（$\nabla$ 为 $TY$ 上的平坦联络），要求
$\nabla$ 给出仿射结构、局部仿射坐标下 $g_{ij}=\p^2K/\p x_i\p x_j$、并满足
Monge--Ampère 方程；去掉最后的方程即得 AK-流形（见 §3.2 与 \cite{CY}）。
本节只用到「度量局部是 Hessian」这一性质，它保证 $K_i-K_j$ 是局部仿射函数。
\end{yizhu}

今后我们假定 $[g]$ 是整的。于是由 GAGA（见 \cite{Be} 命题 3.14），$Y^{an}$ 上解析
凝聚层的范畴等价于相应的射影代数簇 $\YY$ 上代数凝聚层的范畴。

层 $\OO_Y$ 容许一个由软微分分次代数层构成的消解 $\widehat{\Omega}_Y^*$。局部地，
对小开集 $U\subset Y$，$\widehat{\Omega}_Y^*$ 的截面由和
\[
  \alpha=\sum_{i_1,\dots,i_n}c_{i_1\dots i_n}z_1^{i_1}\cdots z_n^{i_n},
  \qquad\text{其中}\quad
  c_{i_1\dots i_n}=\sum_j c_{j,i_1\dots i_n}e^{-\lam_{j,i_1\dots i_n}/\eps}
\]
给出，这里 $c_{j,i_1\dots i_n}\in\Omega^*(U)$，并满足与层 $\OO_Y$ 相同的收敛条件。
微分由作用在系数 $c_{j,i_1\dots i_n}$ 上的 de Rham 微分给出。

我们如下定义微分分次范畴 $C(Y)$。对象是有限秩局部自由 $\OO_Y$-模的有限复形。
对任意两个这样的复形 $E_1$ 与 $E_2$，态射空间定义为
\[
  \Hom_{C(Y)}(E_1,E_2)
  =\Gamma\bigl(Y,\Hom_{\OO_Y}(E_1,E_2)\widehat{\otimes}_{\OO_Y}\widehat{\Omega}_Y^*\bigr),
\]
其中右端使用完备张量积。态射空间上的微分与分次由 $E_1$、$E_2$、$\widehat{\Omega}_Y^*$
上的相应结构诱导。我们把 $C(Y)$ 视为一个 $\Ainf$-预范畴，其中所有对象序列都横截，
且除 $m_1$ 与 $m_2$ 外没有高阶复合。

对域上给定的射影代数簇 $V$，可以典范地定义一个 $\Ainf$-范畴的等价类
$D^b_{\Ainf}(V)$。它由 $V$ 上凝聚层的有界导出范畴增强而得：对象与凝聚层导出范畴
相同；为定义两个对象之间的态射空间，先把它们替换为任意选定的、由局部自由层给出的
零调消解（例如 Godement 消解），再取消解之间态射空间在层复形范畴中的整体截面。
这样就得到一个微分分次范畴。在复数域上的射影簇情形，还有一种用全纯向量丛复形与
Dolbeault 形式给出的替代构造。
消解的不同选取导致 $\Ainf$-等价的范畴。我们把这类范畴的（$\Ainf$-等价）类记为
$D^b_{\Ainf}(V)$。

由定义，$C(Y)$ 的态射空间是相应 $\OO_Y$-模层空间的消解；于是由 \cite{Be} 中的
GAGA 定理，下述结果成立。

\begin{prp}
范畴 $C(Y)$ 与 $D^b_{\Ainf}(\YY)$ 是 $\Ainf$-等价的，其中 $\YY$ 是对应于指派给
$Y$ 的解析空间 $Y^{an}$ 的射影代数簇。
\end{prp}

% [原文 p48–p52]
